\documentclass[11pt]{article}

\usepackage[margin=0.85in,top=0.6in,bottom=0.6in]{geometry}
\usepackage{newtxtext,newtxmath}
\usepackage[T1]{fontenc}
\usepackage{enumitem}
\usepackage{parskip}
\usepackage{titlesec}
\usepackage[hidelinks]{hyperref}

\setlength{\parskip}{4pt}
\titlespacing*{\section}{0pt}{9pt}{3pt}
\titleformat{\section}{\large\bfseries}{\thesection}{0.6em}{}
\setlist[itemize]{leftmargin=1.6em, itemsep=2pt, topsep=2pt}

\pagenumbering{gobble}

\begin{document}

\begin{center}
{\LARGE\bfseries Interpretable Respiratory Sound Analysis}\\[2pt]
{\normalsize Project Background}
\end{center}

\vspace{2pt}
\noindent This note explains what the project is, where the recordings come from, and why your
listening judgment is genuinely useful to it.

\section*{What the project is}
Most machine learning models for detecting disease from lung sounds work as a black box: audio
goes in, a diagnosis comes out, and there is no way for a clinician to see or correct the
reasoning in between. This project builds a diagnosis pipeline that instead routes every decision
through a small set of \textbf{named acoustic concepts} --- crackle, wheeze, and related
adventitious sounds --- computed directly from the audio signal using established signal-processing
methods, not learned from opaque features. The disease prediction is only allowed to depend on
these concept values, so a clinician can, in principle, inspect which sounds the model believed it
heard, and correct one if it is wrong.

The immediate technical question is simpler than that final goal: \textbf{do the automated
concept detectors actually measure what they claim to measure?} A detector labeled ``crackle''
is only useful if it responds to crackles specifically, not to breath sounds in general. That is
what this listening task is designed to check.

\section*{The dataset}
The audio comes from the \textbf{ICBHI 2017 Respiratory Sound Database}, a public research corpus
of 920 recordings from 126 patients, collected with digital stethoscopes across multiple clinical
sites. Each recording is segmented into individual respiratory cycles, and each cycle is already
labeled by the dataset's original creators for the \emph{presence} of crackles and wheezes --- but
not for finer distinctions such as fine versus coarse crackle, or wheeze pitch. All recordings are
anonymized archival research data; there is no patient-identifying information, clinical history,
or imaging attached to any clip.

\section*{Why we need your review}
The dataset's own labels only go so far: they say whether a crackle is present, not what kind. Our
automated detectors compute several finer-grained acoustic properties, but without an independent
human judgment there is no way to know whether those finer measurements are trustworthy. Your
listening task provides exactly that independent check --- a blind, expert judgment on a sample of
clips, made without seeing what our detectors concluded, so that your answers can be compared
fairly against both the automated output and the dataset's own labels. This also lets us tell
apart two different explanations for any disagreement we see: that our detectors need work, or
that the dataset's original labels are themselves imperfect --- both are useful things to know, and
both are publishable findings in their own right.

\section*{Publication scope}
This work is being carried out as a university machine learning research project, with the
explicit goal of submitting the results to a peer-reviewed biomedical engineering journal. Your
listening judgments, if used, would be described in the resulting paper as an expert validation
step, and you would be acknowledged by name (with your permission) as the clinician who performed
it. No patient data, institutional affiliation, or clinical opinion beyond your acoustic
judgments on these anonymized clips would be attributed to you.

\vspace{6pt}
\noindent Thank you again for lending your expertise to this --- it directly strengthens the
scientific credibility of the project.

\end{document}
